\documentclass[../main.tex]{subfiles}
\begin{document}

% Kiến trúc phần mềm, công cụ, kiểm thử và tái lập, chuyển từ Chương 5. Nguồn: repo SUMO (pyproject.toml, src/thesis_sumo/**,
% tests/**, scripts/setup_env.py, CLAUDE.md).

\section{Mạng sinh từ OpenStreetMap}

Hình~\ref{fig:3.net} là mạng N7 do script sinh từ OpenStreetMap với cấu hình hiện trường của nút nghiên cứu, hiển thị trong \texttt{sumo-gui} theo hướng bản đồ. Sơ đồ phân làn ở Hình~\ref{fig:3.1} vẽ từ cùng tệp mạng.

\begin{figure}[!htbp]
\centering
\includegraphics[width=0.95\textwidth]{Figure/net_site_gui.png}
\caption{Mạng N7 trong \texttt{sumo-gui}, hướng bản đồ, khung rộng khoảng 200~m. Lê Văn Lương chạy từ góc dưới trái (LVL\_W, phía Tố Hữu) lên góc trên phải (LVL\_E, phía Láng Hạ); HMG ở góc trên trái, NT ở góc dưới phải, cả hai có dải phân cách ở đoạn cuối. Màu xanh nhạt: đoạn 30~m của làn rẽ phải liên tục; ô hồng: vòng từ đếm xe của bộ tính chỉ số; vạch đỏ và xanh ở mép nút: trạng thái đèn của từng liên kết.}
\label{fig:3.net}
\end{figure}

\section{Kiến trúc phần mềm}
\label{section:architecture}

Hệ thống là một đường ống gồm các bước tách rời; mỗi bước đọc tệp đầu vào và ghi tệp đầu ra dạng văn bản (YAML, JSON, XML, CSV). Cách tổ chức này đáp ứng ba yêu cầu. Yêu cầu thứ nhất là tái lập: mạng sinh lại từ dữ liệu gốc bằng script và tệp vá, và mọi lần chạy dùng seed cố định và ghi \texttt{manifest.json}. Yêu cầu thứ hai là thay được nguồn đặc tả: B0, B1, Optuna và LLM đều chỉ tạo một đặc tả JSON cùng lược đồ. Yêu cầu thứ ba là dùng chung hàm chỉ số: đầu ra SUMO được chuyển về lược đồ của dữ liệu quan sát.

Thời gian chạy chia hai giai đoạn. Giai đoạn ngoại tuyến gồm sinh hoặc soạn đặc tả, kiểm tra hợp lệ và biên dịch; LLM chỉ tham gia ở đây và chỉ tạo tệp JSON. Giai đoạn trực tuyến là một lần chạy SUMO: SUMO nạp tham số tĩnh từ \texttt{vtypes.add.xml} khi khởi động, còn bộ điều khiển TraCI đọc \texttt{traci\_rules.json} và đánh giá luật ở mỗi bước. Một trình chạy duy nhất điều phối mỗi lần chạy: xác định đầu vào theo quy ước tên, khởi động SUMO qua \texttt{libsumo} (chạy hàng loạt) hoặc \texttt{traci} (khi cần \texttt{sumo-gui}), gắn bộ điều khiển nếu đặc tả có luật, ghi \texttt{manifest.json} và chuyển đầu ra sang lược đồ quan sát. Các seed chạy song song, mỗi tiến trình một phiên SUMO. Bảng~\ref{table:4.2} mô tả chín bước của luồng dữ liệu.

\begin{table}[!htbp]
\centering
\small
\caption{Luồng dữ liệu giữa các gói.} \label{table:4.2}
\begin{tabular}{|p{0.6cm}|>{\raggedright\arraybackslash}p{3.4cm}|>{\raggedright\arraybackslash}p{3.6cm}|>{\raggedright\arraybackslash}p{4.9cm}|}
\hline
\textbf{\#} & \textbf{Bước (gói)} & \textbf{Đầu vào} & \textbf{Đầu ra} \\ \hline
1 & Dựng mạng (\texttt{net}) & Bản xuất OSM, tệp cấu hình hiện trường & Mạng cơ sở, mạng hiện trường, tệp vá, bảng cạnh, sơ đồ mạng \\ \hline
2 & Dữ liệu mục tiêu (\texttt{proxy}) & Bảng giá trị mục tiêu, cấu hình hiện trường & Các bảng CSV theo lược đồ quan sát của từng ngày \\ \hline
3 & Nhu cầu (\texttt{demand}) & Mạng, số đếm rẽ 15 phút & Tệp tuyến của từng ngày \\ \hline
4 & Kiểm tra hợp lệ (\texttt{rules}) & Đặc tả JSON, miền tham số, cấu hình hiện trường, mạng & Danh sách lỗi có mã \\ \hline
5 & Biên dịch (\texttt{rules}) & Đặc tả hợp lệ & \texttt{vtypes.add.xml}, \texttt{traci\_\allowbreak rules.json} (nếu có luật) \\ \hline
6 & Chạy mô phỏng (\texttt{sim}, \texttt{runtime}) & Mạng, tệp tuyến, đặc tả đã biên dịch, cấu hình thí nghiệm, seed & Đầu ra SUMO, \texttt{sumo.log}, \texttt{interventions.csv}, \texttt{manifest.json} \\ \hline
7 & Chuyển lược đồ (\texttt{sim}) & Đầu ra SUMO & Các bảng CSV theo lược đồ quan sát \\ \hline
8 & Chỉ số và điểm (\texttt{metrics}, \texttt{evaluate}) & Bảng CSV mô phỏng và quan sát của các ngày hiệu chỉnh & M1--M10, G0, $R$, $S$, tiêu chí chấp nhận \\ \hline
9 & Sinh đặc tả kế tiếp (\texttt{llm}, \texttt{baselines}) & Kết quả đánh giá, ngân sách còn lại & Phản hồi cho LLM hoặc giá trị mục tiêu cho Optuna; đặc tả mới \\ \hline
\end{tabular}
\end{table}

Phụ thuộc giữa các gói đi một chiều. \texttt{config} và \texttt{schema\_observed} là nền. \texttt{rules} và \texttt{runtime} dùng chung một mô-đun xác định làn đi sai dùng được, để bộ kiểm tra và bộ điều khiển hiểu giống nhau. \texttt{proxy} sinh dữ liệu độc lập với đầu ra SUMO, nhưng dùng lại định nghĩa hình học và lược đồ của \texttt{sim}. \texttt{metrics} chỉ đọc các bảng CSV theo lược đồ quan sát nên không phụ thuộc SUMO. Bảng~\ref{table:4.1} nêu trách nhiệm và quyết định thiết kế của từng gói.

\begingroup\small
\begin{longtable}{|>{\raggedright\arraybackslash}p{2.3cm}|>{\raggedright\arraybackslash}p{4.2cm}|>{\raggedright\arraybackslash}p{7.1cm}|}
\caption{Các gói con của \texttt{thesis\_sumo} và quyết định thiết kế.}\label{table:4.1}\\
\hline
\textbf{Gói} & \textbf{Trách nhiệm} & \textbf{Quyết định thiết kế} \\ \hline
\endfirsthead
\hline
\textbf{Gói} & \textbf{Trách nhiệm} & \textbf{Quyết định thiết kế} \\ \hline
\endhead
\texttt{config}, \texttt{schema\_\allowbreak observed} & Đọc cấu hình hiện trường, thí nghiệm, LLM; định nghĩa bảng và cột của lược đồ quan sát & Trường bắt buộc rỗng gây lỗi \texttt{MissingFieldDataError} nêu tên trường, không điền mặc định. Tệp đo thực mang nhãn \texttt{MEASURED}; tệp thay thế mang nhãn \texttt{LITERATURE\_PROXY}, mỗi giá trị kèm nguồn và nhãn bằng chứng. Chương trình không bao giờ tự chuyển từ tệp đo thực sang tệp thay thế \\ \hline
\texttt{net} & Sinh mạng từ OSM, áp tệp vá, xuất bảng cạnh và sơ đồ mạng & Mạng cơ sở qua hai lượt \texttt{netconvert}. Mạng hiện trường chạy lại \texttt{netconvert} tới khi điểm tách cạnh đúng khoảng cách (tối đa 5 lượt), rồi thêm lượt cuối đặt lại điểm chờ của xe rẽ trái trái phép. Việc tách cạnh làm vạch dừng dịch từ $-0{,}70$ đến $+1{,}61$~m tùy nhánh; độ dịch ghi vào một tệp để các bảng quan sát dùng đúng vị trí. Mạng N6 có 28 liên kết đèn: 24 liên kết lòng đường và 4 liên kết vỉa hè \\ \hline
\texttt{demand} & Sinh tuyến ứng viên và tệp tuyến từ số đếm rẽ 15 phút & Dùng công cụ \texttt{routeSampler} của SUMO \\ \hline
\texttt{proxy} & Sinh năm ngày dữ liệu thay thế PROXY\_D1--PROXY\_D5 từ bảng giá trị mục tiêu & Mỗi ngày gồm mười bảng CSV, một manifest và tệp tuyến. Không đọc đầu ra SUMO. Ba bảng của nhóm tìm khoảng trống dùng luồng số ngẫu nhiên riêng để các bảng cũ không đổi \\ \hline
\texttt{rules} & Lược đồ đặc tả, bộ kiểm tra hợp lệ, bộ biên dịch & Bộ kiểm tra trả cả danh sách lỗi, nội dung nêu rõ khoảng hoặc tên được phép để đưa nguyên văn vào phản hồi. Bộ biên dịch gọi bộ kiểm tra trước; đặc tả sai gây \texttt{InvalidSpecError} và không ghi tệp. Mọi tham số, kể cả bằng mặc định, được ghi tường minh. Tệp sinh ra khai báo XSD của SUMO, vì nếu không khai báo, thuộc tính sai tên bị bỏ qua không cảnh báo. Nguyên mẫu \texttt{custom1} (\texttt{custom2}) nhận trước giá trị mặc định của xe máy (ô tô), vì vClass tùy biến không có mặc định này \\ \hline
\texttt{runtime} & Bộ điều khiển TraCI và sáu luật & Mỗi luật là lớp con của \texttt{Rule}, đăng ký trong \texttt{RULE\_REGISTRY}. Mỗi luật có bộ sinh số ngẫu nhiên riêng, khởi tạo từ seed mô phỏng, chỉ số và loại luật; xe duyệt theo thứ tự mã, nên \texttt{libsumo} và \texttt{traci} cho cùng chuỗi rút và cùng nhật ký. Mọi lệnh đi qua lớp bao \texttt{Commands}: mỗi xe một can thiệp, chỉ xóa đúng bit được nêu của speed mode, kiểm tra lại giá trị SUMO báo về \\ \hline
\texttt{sim} & Ghi cấu hình SUMO, chạy, ghi manifest, chuyển đầu ra & \texttt{run.sumocfg} do chính SUMO ghi, nên tùy chọn sai tên bị từ chối. Luôn bật \texttt{collision.check-junctions}. Ba bảng của nhóm tìm khoảng trống chỉ suy từ quỹ đạo, detector và tuyến, không đọc nhật ký can thiệp \\ \hline
\texttt{metrics} & Chỉ số M1--M10, G0, hàm khoảng cách, điểm tổng & Mỗi chỉ số một mô-đun, cộng một mô-đun dùng chung (gán chu kỳ đèn, xác định xe rẽ phải liên tục và các hướng có mũi tên riêng để loại khỏi chỉ số vi phạm). Hàm \texttt{compare} trả $S$, $S_{\text{sourced}}$, $S_{\text{assumed}}$, khoảng cách và $R$ theo khóa \\ \hline
\texttt{evaluate} & Ngân sách theo chiến dịch, đánh giá một đặc tả, chọn đặc tả tốt nhất & \texttt{RunBudget} giữ ngân sách trong \texttt{budget.json}, đọc lại từ đĩa, trừ dưới khóa tệp. Đặc tả không hợp lệ không tốn lần chạy; mã trùng trong một chiến dịch bị từ chối; ngày giữ riêng bị chặn \\ \hline
\texttt{llm} & Gọi LLM, prompt có phiên bản, phản hồi, vòng lặp P1/P2 & Hai backend: \texttt{openai} (yêu cầu khai báo nhiệt độ và mức lượng tử hóa) và \texttt{file} (ghi yêu cầu ra tệp và chờ tệp phản hồi, khi máy chủ LLM ở máy khác). Phát lại nhật ký có kiểm định danh LLM. \texttt{run\_loop} lưu trạng thái chiến dịch và từ chối chạy tiếp nếu trạng thái đổi \\ \hline
\texttt{baselines} & Optuna cho B2 & Nghiên cứu Optuna lưu trong tệp nhật ký, chạy tiếp không trừ lại ngân sách, từ chối nếu ngày, seed hoặc cấu trúc khác. Đặc tả tốt nhất chọn bằng cùng \texttt{select\_best} với vòng LLM \\ \hline
\end{longtable}
\endgroup

Lớp cơ sở \texttt{Rule} giữ tham số luật, tập mã vType áp dụng, thông tin nút, lớp bao lệnh, bộ sinh số ngẫu nhiên và một từ điển các can thiệp đang diễn ra. Phương thức trừu tượng \texttt{on\_step} nhận ảnh chụp trạng thái của bước và trả về danh sách can thiệp đã có kết quả trong bước đó. Lớp cơ sở cung cấp năm thao tác dùng chung: kiểm tra xe thuộc phạm vi luật và đang rảnh; bắt đầu can thiệp; hoàn tác ghi đè; kết thúc can thiệp kèm kết quả và lý do; đóng mọi can thiệp còn dang dở khi lần chạy kết thúc. Khi SUMO báo speed mode khác giá trị vừa đặt, lớp bao lệnh báo \texttt{ModeNotRestoredError}.

\section{Thư viện và công cụ}
\label{section:tools}

Bảng~\ref{table:4.3} liệt kê công cụ và thư viện. Cột ``Yêu cầu'' là ràng buộc tối thiểu khai báo trong cấu hình gói; cột ``Đã cài'' là phiên bản trong môi trường ảo của máy chính.

\begin{table}[!htbp]
\centering
\small
\caption{Thư viện và công cụ sử dụng.} \label{table:4.3}
\begin{tabular}{|>{\raggedright\arraybackslash}p{4.6cm}|>{\raggedright\arraybackslash}p{4.4cm}|p{2.0cm}|p{2.0cm}|}
\hline
\textbf{Mục đích} & \textbf{Công cụ} & \textbf{Yêu cầu} & \textbf{Đã cài} \\ \hline
Mô phỏng vi mô, dựng mạng, sinh tuyến & SUMO (\texttt{sumo}, \texttt{sumo-gui}, \texttt{netconvert}, \texttt{routeSampler})~\cite{lopez2018sumo} & 1.27.0 & 1.27.0 \\ \hline
Điều khiển SUMO từ Python, chạy hàng loạt & \texttt{libsumo}, \texttt{sumo-data} & 1.27.0 & 1.27.0 \\ \hline
Điều khiển SUMO qua giao diện đồ họa, đọc mạng & \texttt{traci}, \texttt{sumolib} (lấy từ thư mục cài đặt SUMO) & như SUMO & 1.27.0 \\ \hline
Ngôn ngữ lập trình & Python (môi trường ảo) & $\geq$ 3.10 & 3.12.10 \\ \hline
Đọc cấu hình & \texttt{pyyaml} & $\geq$ 6.0 & 6.0.3 \\ \hline
Tính toán số & \texttt{numpy} & $\geq$ 1.26 & 2.5.4 \\ \hline
Xử lý bảng & \texttt{pandas} & $\geq$ 2.1 & 3.1.0 \\ \hline
Thống kê & \texttt{scipy} & $\geq$ 1.11 & 1.18.1 \\ \hline
Đọc và ghi XML & \texttt{lxml} & $\geq$ 5.0 & 6.1.3 \\ \hline
Vẽ hình & \texttt{matplotlib} & $\geq$ 3.8 & 3.11.2 \\ \hline
Kiểm tra đặc tả theo lược đồ JSON & \texttt{jsonschema} & $\geq$ 4.20 & 4.26.0 \\ \hline
Chuyển tọa độ kinh vĩ độ sang tọa độ mạng & \texttt{pyproj} & $\geq$ 3.6 & 3.8.0 \\ \hline
Kiểm thử tự động & \texttt{pytest} & $\geq$ 8.0 & 9.1.1 \\ \hline
Gọi LLM qua giao diện tương thích OpenAI & \texttt{openai} & $\geq$ 1.30 & 3.40.0 \\ \hline
Tối ưu tham số B2~\cite{akiba2019optuna} & \texttt{optuna} & $\geq$ 3.5 & 5.0.0 \\ \hline
Chạy LLM cục bộ & Ollama; mô hình \texttt{qwen3.5:9b}, Q4\_K\_M & -- & Ollama 0.40.11 \\ \hline
Quản lý phiên bản mã & Git (kèm Git LFS cho tệp lớn) & -- & 2.49.0 \\ \hline
\end{tabular}
\end{table}

\texttt{traci} và \texttt{sumolib} không cài từ PyPI mà nạp từ thư mục công cụ của bản cài SUMO, để khớp với tệp thực thi. \texttt{libsumo} và \texttt{sumo-data} cài không kèm phụ thuộc, vì phụ thuộc khai báo của chúng kéo theo \texttt{traci} và \texttt{sumolib} bản khác từ PyPI và che bản đi kèm SUMO. Trên máy chính, nếu \texttt{optuna} được nạp trước \texttt{libsumo} thì \texttt{libsumo} không nạp được, vì tệp \texttt{sqlite3.dll} đi kèm Python xung đột với tệp cùng tên đi kèm \texttt{libsumo}; script Optuna vì vậy nạp \texttt{libsumo} trước. Tệp nhật ký nghiên cứu của Optuna khóa bằng \texttt{JournalFileOpenLock}, vì cách khóa mặc định tạo liên kết tượng trưng, mà tài khoản Windows thường không có quyền tạo.

\section{Các bộ kiểm thử}

Bảng~\ref{table:4.5} tóm tắt các bộ kiểm thử tự động; kết quả nằm ở Mục~\ref{section:testing}.

\begin{table}[!htbp]
\centering
\small
\caption{Các bộ kiểm thử tự động.} \label{table:4.5}
\begin{tabular}{|>{\raggedright\arraybackslash}p{3.6cm}|>{\raggedright\arraybackslash}p{9.4cm}|}
\hline
\textbf{Bộ kiểm thử} & \textbf{Nội dung kiểm tra} \\ \hline
\texttt{test\_env} & \texttt{traci} và \texttt{sumolib} nạp từ bản cài SUMO và khớp phiên bản tệp thực thi \\ \hline
\texttt{test\_config} & Trường bắt buộc rỗng gây lỗi nêu tên trường; cấu hình đo thực không bao giờ tự lấy giá trị thay thế \\ \hline
\texttt{test\_net} & Dựng mạng bằng một lệnh; tệp vá, đèn, bảng cạnh; dựng lại cho kết quả giống từng byte; làn nội bộ có bề rộng bằng làn ra tương ứng \\ \hline
\texttt{test\_demand} & vType B0, tuyến ứng viên, sinh tuyến từ số đếm rẽ, kiểm tra dữ liệu số đếm \\ \hline
\texttt{test\_sumocfg} & Tùy chọn SUMO lấy từ cấu hình thí nghiệm, tệp đầu ra, chế độ SSM, chạy thật ngắn \\ \hline
\texttt{test\_runner} & Trình chạy, trường của manifest, chuyển sang lược đồ quan sát, tiêu chí chấp nhận \\ \hline
\texttt{test\_outputs\_\allowbreak synthetic} & Chuyển lược đồ trên đầu ra SUMO viết tay cho mạng chữ thập nhỏ, kết quả kỳ vọng tính tay \\ \hline
\texttt{test\_proxy} & Tính tất định của bộ sinh dữ liệu thay thế, lược đồ, nhất quán nội tại, độ khớp với giá trị mục tiêu \\ \hline
\texttt{test\_validator} & Miền tham số, tham số luật, mã lỗi và nội dung thông báo lỗi; lựa chọn cắt về \\ \hline
\texttt{test\_compiler} & Biên dịch B1 khớp từng byte với bản chụp; mọi thuộc tính nằm trong miền; SUMO nạp được với kiểm tra XSD; đặc tả B0 tương đương tệp vType viết tay \\ \hline
\texttt{test\_runtime} & Mỗi luật một kịch bản trên mạng nhỏ; lệnh bị SUMO bỏ qua được ghi là không thực thi; khôi phục speed mode; giữ xe giữa nút không gây va chạm trong nút; cấm \texttt{ignoreIDs} trong mã \\ \hline
\texttt{test\_backends} & \texttt{libsumo} và \texttt{traci} cho quỹ đạo giống hệt nhau với cùng kịch bản và seed \\ \hline
\texttt{test\_metrics} & M1--M10, G0, hàm khoảng cách trên bảng dựng tay có đáp án; khoảng cách khi một phía rỗng không bao giờ nhỏ hơn khoảng cách khi sai \\ \hline
\texttt{test\_evaluate} & Ngân sách không vượt khi chạy lại; đặc tả hỏng không tốn lần chạy; mã trùng bị từ chối; chặn ngày giữ riêng; quy tắc chọn đặc tả tốt nhất \\ \hline
\texttt{test\_llm} & Lược đồ cho LLM, khóa cấu trúc, phát lại nhật ký có kiểm định danh LLM, chặn rò rỉ dữ liệu ngày giữ riêng \\ \hline
\texttt{test\_loop} & Chạy tiếp đúng sau khi bị ngắt mà không trừ lại ngân sách; mẫu không hợp lệ không tốn lần chạy; khóa cấu trúc của P1; chạy thêm vòng khi còn ngân sách \\ \hline
\texttt{test\_optuna} & Không gian tìm kiếm theo cấu trúc; mọi thử nghiệm ngẫu nhiên đều hợp lệ; chạy tiếp không trừ lại ngân sách \\ \hline
\end{tabular}
\end{table}

\section{Triển khai và tái lập}
\label{section:deploy}

\textbf{Môi trường.} Máy chính chạy Windows 11 Pro với CPU Intel Core i5-13400F (10 nhân, 16 luồng), 31,8~GB RAM và GPU NVIDIA GeForce RTX 4060 8~GB; SUMO 1.27.0 cài sẵn, biến \texttt{SUMO\_HOME} trỏ tới thư mục cài đặt, và LLM cục bộ chạy trên máy này. Một máy Linux phụ, với SUMO 1.27.0 cài dưới dạng gói Python \texttt{eclipse-sumo} trong môi trường ảo, được dùng để kiểm tra tái lập: tám lần chạy B0 trên máy chính được chạy lại trên máy này và cho đúng từng số va chạm và teleport. Môi trường Python là môi trường ảo Python 3.12, dựng bằng cùng một script trên cả hai máy.

\begin{sloppypar}
\textbf{Lệnh tái lập.} Các lệnh mô-đun Python theo thứ tự của Bảng~\ref{table:4.2}: (i) \texttt{thesis\_sumo.net.build} với tệp cấu hình hiện trường \texttt{config/\allowbreak site.yaml}: dựng mạng N7; (ii) \texttt{video/\allowbreak make\_video\_days.py}: sinh hai ngày VIDEO\_V2, VIDEO\_V1 và tệp tuyến bằng \texttt{thesis\_sumo.proxy.generate} trên bảng mục tiêu đo từ video (với tệp cấu hình thay thế, cùng mô-đun sinh năm ngày PROXY); (iii) \texttt{thesis\_sumo.rules.compiler}: biên dịch một đặc tả; (iv) \texttt{thesis\_sumo.sim.runner}: chạy mô phỏng, với tham số mã đặc tả, ngày, danh sách seed, tệp cấu hình hiện trường, thời điểm kết thúc, số tiến trình song song và tùy chọn mở \texttt{sumo-gui}; và (v) \texttt{thesis\_sumo.baselines.optuna\_search} (B2) và \texttt{thesis\_sumo.llm.loop} (P1 với chế độ \texttt{params\_only} và cơ sở B1, P2 với chế độ \texttt{full}): chạy chiến dịch với cùng ngân sách, cùng cấu hình hiện trường và cùng mức trần va chạm của B0. Các script của Chương~\ref{chapter:Experiment} trên N7 nằm trong \texttt{experiments/\allowbreak operator/\allowbreak n1126/}: \texttt{run\_campaigns\_video.sh} (chiến dịch), \texttt{eval\_benchmark\_video.py} (benchmark), \texttt{summarise.py} và \texttt{compare\_video.py} (bảng và hình). Dữ liệu mỗi chiến dịch nằm trong \texttt{experiments/\allowbreak campaigns/\allowbreak <chiến dịch>/}; nhật ký LLM nằm trong \texttt{experiments/\allowbreak llm\_logs/\allowbreak <chiến dịch>/}.
\end{sloppypar}

\textbf{Đầu ra của một lần chạy.} Mỗi lần chạy ghi vào \texttt{experiments/\allowbreak runs/\allowbreak <spec\_id>\_\_<day>/\allowbreak seed\_XX/}. Thư mục gồm \texttt{run.sumocfg}; đầu ra SUMO (thông tin chuyến đi, quỹ đạo nén, chuyển làn, thống kê, va chạm, trạng thái đèn, detector); \texttt{sumo.log}; \texttt{interventions.csv} nếu có luật TraCI; \texttt{manifest.json} (mã commit, phiên bản SUMO và backend, giá trị băm của đặc tả, mạng, tệp tuyến và cấu hình hiện trường, cấu hình thí nghiệm, số teleport và va chạm theo loại, tóm tắt can thiệp, số xe chờ vào mạng, cờ \texttt{valid\_strict}); và các bảng CSV theo lược đồ quan sát. Đầu ra thô không đưa vào kho mã; kho mã lưu mã, cấu hình, tệp vá, mạng đã sinh, dữ liệu thay thế, \texttt{manifest.json} và các bảng kết quả nhỏ. Một lần chạy B0 1800~s trên dữ liệu thay thế, khi chạy 3 lần song song, mất 88--463~s.

\end{document}
